\newpage
\section{Projektowanie}		%3
%Napisać z jakich narzędzi będziemy korzystać (kompilator, język programowania), git, biblioteki dodatkowe, itp.
%Opisać szczegółowe ustawienia kompilatora (jeśli są), powiązania z bibliotekami, itp.
%Narysować graf, UML, diagram klas, schemat działania algorytmu
%Jeśli zadanie zakłada przedstawienie jakiegoś narzędzia (np. git, AI) należy opisać sposób jego używania

\hspace{0.60cm}
W niniejszym rozdziale zaprezentowano etap projektowy aplikacji, obejmujący zestawienie użytych narzędzi, konfiguracyjne ustawienia kompilatora, architekturę klas wyrażoną diagramem UML oraz opis procedury zarządzania repozytorium w systemie kontroli wersji.

\subsection{Środowisko programistyczne i konfiguracja kompilatora}

\hspace{0.60cm}
Do realizacji projektu wykorzystano standard języka \textbf{C++17}, zapewniający wsparcie dla kontenerów STL (\texttt{std::vector}, \texttt{std::string}) oraz ulepszoną obsługę strumieni plików (\texttt{std::fstream}). Zespół pracował w środowiskach Microsoft Visual Studio 2022 oraz Visual Studio Code z wykorzystaniem kompilatorów MSVC oraz MinGW (GCC 12.2).

\hspace{0.60cm}
Konfiguracja kompilatora obejmowała następujące flagi budowania projektu:
\begin{itemize}
    \item \texttt{-std=c++17} (lub \texttt{/std:c++17}) -- wymuszenie standardu C++17.
    \item \texttt{-Wall -Wextra} (lub \texttt{/W4}) -- włączenie szczegółowych ostrzeżeń kompilatora w celu wykrycia potencjalnych wycieków pamięci i niezainicjalizowanych wskaźników.
    \item \texttt{-O2} -- flaga optymalizacyjna dla wersji ostatecznej (Release).
\end{itemize}

\subsection{Diagram klas UML i architektura systemu}

\hspace{0.60cm}
Architektura aplikacji opiera się na ścisłym podziale odpowiedzialności. Główna logika drzewa została odizolowana od modułu I/O oraz warstwy interfejsu. Rysunek \ref{rys:uml_diagram} (s. \pageref{rys:uml_diagram}) przedstawia diagram klas UML opisujący powiązania pomiędzy strukturami programu.

%\begin{figure}[!htb]
%	\begin{center}
%		\framebox{\parbox{11cm}{\centering \vspace{2.2cm} [ Diagram klas UML: Node $\leftarrow$ BSTTree $%\leftrightarrow$ FileHandler $\leftarrow$ main ] \vspace{2.2cm}}}
%		\caption{Diagram klas UML aplikacji drzewa BST}
%		\label{rys:uml_diagram}
%	\end{center}
%\end{figure}

\subsection{Przepływ pracy i sposób użycia narzędzia Git}

\hspace{0.60cm}
Praca w zespole dwuosobowym odbywała się z wykorzystaniem modelu \textit{Feature Branch Workflow}. Poniżej opisano sekwencję poleceń używanych podczas tworzenia nowych funkcjonalności i rozwiązywania konfliktów:

\begin{enumerate}
    \item \textbf{Tworzenie gałęzi:} \texttt{git checkout -b feature/nazwa\_funkcji} -- praca na wyizolowanej gałęzi roboczej.
    \item \textbf{Zatwierdzanie etapów:} \texttt{git add .} oraz \texttt{git commit -m "Opis zmian"} -- tworzenie lokalnych punktów kontrolnych.
    \item \textbf{Publikacja i scalanie:} \texttt{git push origin feature/...} oraz wywołanie \texttt{git merge} do gałęzi \texttt{main}.
    \item \textbf{Rozwiązywanie konfliktów:} W przypadku wystąpienia konfliktu scalania używano wbudowanego edytora \textit{Merge Editor} w środowisku VS Code/Visual Studio do ręcznego wybrania poprawnych linii kodu.
\end{enumerate}